%% 94_appendix_search.tex -- split from main.tex by paper-planner (2026-10-01). Owner group: G4.
%% Camera-ready (G4, 2026-10-01): App. E, search cost and selector controls (yHj7-1, yHj7-2,
%% yHj7-4, ws2F-6, INT-13). Macros: numbers_cr_search.tex (\Ncr*), numbers.tex (\Nscope*),
%% numbers_v14.tex (\Nmenc*).
\section{Search: cost and selector controls}
\label{app:search}
Unless stated otherwise, every search result uses the \NmencRuns\ runs of the \Nbudgetc\ budget and
\NmencItems\ evaluation items per run. Every candidate is drawn at $T{=}8$ (64 NFE), and candidates
differ only in their sampling noise. Best-of-$K$ search therefore uses NFE $=64K$ and is
matched to refinement at $T{=}8K$. A gain is a difference of per-run means, with a bootstrap
over runs (\Nbootreps\ replicates). A win rate is the share of items on which search scores
higher, with a bootstrap over items.

\paragraph{Selection cost.}
Refinement decodes its output once and runs no selector. Search decodes all $K$ candidates
with Mimi and embeds each of them, and the prompt once, with WavLM-L. We count the
generator's FLOPs analytically and check the count against PyTorch's FLOP counter. We count
Mimi and WavLM-L with the same counter on randomly initialised copies, on one item of median
length, and scale to the mean item length. One forward pass costs \NcrGenGflops\ GFLOPs for
the generator (range over the five configurations), \NcrMimiGflops\ for a Mimi decode and
\NcrWavlmGflops\ for WavLM-L. The generator's input includes the phoneme tokens, whose
number we estimate at \NcrPhonPerChar\ phones per character (espeak-ng). Leaving them out
makes the generator cheaper and gives the upper bound in Table~\ref{tab:cost}. The count
leaves out the ECAPA-TDNN-small head on WavLM-L and the 24-to-16\,kHz resampling. The
wall-time ratio sums single forward passes of each network, timed on one item on CPU for
configuration C3; it is not a timed run of the whole sampler. The ECAPA scorer and the ASR
are evaluation-only and are not part of any deployed cost.

\begin{table}[h]
  \centering\small
  \caption{Selection cost of best-of-$K$ search, over refinement at the same NFE. Ranges span
  the five \Nbudgetc\ configurations. The upper bound leaves the phoneme tokens out of the
  generator count. Wall time is relative to refinement.}
  \label{tab:cost}
  \begin{tabular}{rrcccc}
    \toprule
    $K$ & NFE & extra generator passes & extra compute (\%) & upper bound (\%)
      & CPU wall time \\
    \midrule
    2 & 128 & \NcrCostTwonfe   & \NcrCostTwopct   & \NcrCostTwobound   & \NcrWallTworatio \\
    4 & 256 & \NcrCostFournfe  & \NcrCostFourpct  & \NcrCostFourbound  & \NcrWallFourratio \\
    8 & 512 & \NcrCostEightnfe & \NcrCostEightpct & \NcrCostEightbound & \NcrWallEightratio \\
    \bottomrule
  \end{tabular}
\end{table}

\paragraph{Equal total compute.}
We also compare search with the cheapest refinement arm whose total compute, counted as above,
is at least that of search. For best-of-2 this arm is $T{=}\NcrEqualTwoT$, and search gains
\NcrEqualTwod\ in ECAPA similarity (CI \NcrEqualTwoci). For best-of-4 it is
$T{=}\NcrEqualFourT$, and search gains \NcrEqualFourd\ (CI \NcrEqualFourci). No refinement
arm we ran costs as much as best-of-8. These gains are larger than the matched-NFE ones,
because refinement's ECAPA similarity falls past $T{=}16$ (its WavLM-L SIM-o rises slightly;
Table~\ref{tab:gapclosed}). Refinement at $T{=}8$ and at $T{=}16$
differ by $\NcrRefineEightMinusSixteen$ (CI $\NcrRefineEightMinusSixteenCI$). So no refinement
arm from $T{=}8$ to $T{=}64$ reaches the similarity of best-of-2.

\paragraph{Selector and scorer.}
Only WavLM-L and ECAPA scored all eight candidates of every item, so only they can act as
selectors. Table~\ref{tab:xsel} crosses the two. The default (WavLM-L selects, ECAPA scores)
and the swapped control (ECAPA selects, WavLM-L scores) both beat refinement at every tier,
and the win-rate intervals of both lie above one half. The swapped gains are smaller:
\NcrXEWLod, \NcrXEWMidd\ and \NcrXEWHid, against \NcrMatchLod, \NcrMatchMidd\ and
\NcrMatchHid. Rows in which one encoder both selects and scores are circular; we show them
for comparison and draw no conclusion from them. The other three encoders of Table~\ref{tab:menc} (x-vector, GE2E and
WavLM-base+) scored only the WavLM-L-selected output, so we could not test them as
selectors.

\begin{table}[h]
  \centering\small
  \caption{Search against matched-NFE refinement, for each selector and scorer
  (\Nbudgetc\ budget). Each cell gives the gain, the win rate in parentheses and the 95\% CI
  of the gain and of the win rate. The first two rows repeat Table~\ref{tab:search} and add
  win rates. Last row: share of items on which the two selectors pick the same candidate.}
  \label{tab:xsel}
  \setlength{\tabcolsep}{2.5pt}
  \begin{tabular}{lccc}
    \toprule
    selector $\to$ scorer & NFE 128 & NFE 256 & NFE 512 \\
    & $K{=}2$ vs $T{=}16$ & $K{=}4$ vs $T{=}32$ & $K{=}8$ vs $T{=}64$ \\
    \midrule
    WavLM-L $\to$ ECAPA
      & \NcrXWELod\ (\NcrXWELowin\%) & \NcrXWEMidd\ (\NcrXWEMidwin\%) & \NcrXWEHid\ (\NcrXWEHiwin\%) \\
    \quad\footnotesize default
      & {\footnotesize\NcrXWELoci; \NcrXWELowinci} & {\footnotesize\NcrXWEMidci; \NcrXWEMidwinci} & {\footnotesize\NcrXWEHici; \NcrXWEHiwinci} \\
    ECAPA $\to$ WavLM-L
      & \NcrXEWLod\ (\NcrXEWLowin\%) & \NcrXEWMidd\ (\NcrXEWMidwin\%) & \NcrXEWHid\ (\NcrXEWHiwin\%) \\
    \quad\footnotesize swapped
      & {\footnotesize\NcrXEWLoci; \NcrXEWLowinci} & {\footnotesize\NcrXEWMidci; \NcrXEWMidwinci} & {\footnotesize\NcrXEWHici; \NcrXEWHiwinci} \\
    WavLM-L $\to$ WavLM-L
      & \NcrXWWLod\ (\NcrXWWLowin\%) & \NcrXWWMidd\ (\NcrXWWMidwin\%) & \NcrXWWHid\ (\NcrXWWHiwin\%) \\
    \quad\footnotesize circular
      & {\footnotesize\NcrXWWLoci; \NcrXWWLowinci} & {\footnotesize\NcrXWWMidci; \NcrXWWMidwinci} & {\footnotesize\NcrXWWHici; \NcrXWWHiwinci} \\
    ECAPA $\to$ ECAPA
      & \NcrXEELod\ (\NcrXEELowin\%) & \NcrXEEMidd\ (\NcrXEEMidwin\%) & \NcrXEEHid\ (\NcrXEEHiwin\%) \\
    \quad\footnotesize circular
      & {\footnotesize\NcrXEELoci; \NcrXEELowinci} & {\footnotesize\NcrXEEMidci; \NcrXEEMidwinci} & {\footnotesize\NcrXEEHici; \NcrXEEHiwinci} \\
    \midrule
    same pick (chance $1/K$)
      & \NcrAgreeLo\% (\NcrChanceLo\%) & \NcrAgreeMid\% (\NcrChanceMid\%) & \NcrAgreeHi\% (\NcrChanceHi\%) \\
    \bottomrule
  \end{tabular}
\end{table}

\paragraph{Replications.}
We repeated the NFE 512 comparison on the \NmaxNfour\,M budget (five configurations, two
seeds each) and on three configurations retrained for 90k steps ($3\times$ the schedule, one
seed each). With WavLM-L selecting and ECAPA scoring, search gains \NscopeDd\ (CI \NscopeDci,
wins on \NscopeDwin\% of items) on the \NmaxNfour\,M budget and \NscopeNd\ (CI \NscopeNci,
\NscopeNwin\%) after retraining. With ECAPA selecting and WavLM-L scoring, it gains
\NcrXEWDd\ (CI \NcrXEWDci, \NcrXEWDwin\%) and \NcrXEWNd\ (CI \NcrXEWNci, \NcrXEWNwin\%). At
NFE 128 the retrained runs gain \NcrXWENLod\ (CI \NcrXWENLoci) with WavLM-L selecting and
\NcrXEWNLod\ (CI $\NcrXEWNLoci$) with ECAPA selecting. With \NcrXNRuns\ runs the second
interval includes zero, so at this tier the retrained runs do not confirm the swapped
control.

\paragraph{Other speaker encoders.}
Table~\ref{tab:menc} rescores the NFE 512 comparison (best-of-8 against $T{=}64$, WavLM-L
selecting) with five speaker encoders. Their cosines sit on different scales, so we compare
win rates: the share of items on which an encoder prefers the search output. A win rate does
not change under a monotone rescaling of an encoder's cosine. All four
encoders other than the selector prefer the search output on \NmencLo--\NmencHi\% of items,
and every interval lies above one half. We did not repeat this rescoring at NFE 128 or 256.

\begin{table}[h]
  \centering\small
  \caption{Search against refinement at NFE 512 (\NmencRuns\ runs, \NmencItems\ items each),
  rescored by ECAPA \citep{ecapa,speechbrain}, x-vector \citep{xvector} (both trained on VoxCeleb
  \citep{voxceleb1,voxceleb2}), GE2E \citep{ge2e}, WavLM-base+ and WavLM-L \citep{wavlm,seedtts}.}
  \label{tab:menc}
  % GENERATED by src/paper_v14.py
\begin{tabular}{llccc}
  \toprule
  encoder & family / role & win rate & 95\% CI & raw $\Delta$ \\
  \midrule
    ECAPA & ECAPA-TDNN & 72.3\% & [70.6, 73.8] & +0.0444 \\
    x-vector & TDNN & 79.0\% & [77.5, 80.5] & +0.0075 \\
    GE2E d-vector & LSTM d-vector & 69.1\% & [67.5, 70.8] & +0.0164 \\
    WavLM-base+ & WavLM, other scale & 64.6\% & [62.9, 66.3] & +0.0173 \\
    WavLM-L$^{\ast}$ & WavLM / \emph{selector} & 84.9\% & [83.6, 86.1] & +0.0803 \\
    \midrule
    \multicolumn{5}{l}{\footnotesize $^{\ast}$selects the candidates, so this row is circular and is not counted as agreement.} \\
  \bottomrule
\end{tabular}

\end{table}
